% Paper 1 (Outline A / thesis-adjusted Option 1)
% MDMP Staff Planning Assistant: open doctrine LoRA adaptation
%
% Thesis (locked; paraphrase in body, restate once in Sec.~5.3):
% A publishable MDMP coaching assistant can be built from open doctrine alone,
% without proprietary algorithms or customer data, and still substantially
% outperform a strong general-purpose base model on held-out staff-planning questions.
%
% Build: pdflatex paper1-open-mdmp-lora && bibtex paper1-open-mdmp-lora && pdflatex ... && pdflatex ...

\documentclass[11pt]{article}

\usepackage[margin=1in]{geometry}
\usepackage{times}
\usepackage{booktabs}
\usepackage{graphicx}
\usepackage{float}
\usepackage{flafter}
\usepackage{url}
\usepackage{xurl}
\usepackage{hyperref}
\usepackage{natbib}
\usepackage{xcolor}
\usepackage{enumitem}

\title{Open Doctrine, Open Weights:\\
Fine-Tuning a Publishable MDMP Staff Planning Assistant with LoRA}
\author{Dr.\ William J.\ L.\ Adams\\
Decision Lens Inc.}
\date{}

\begin{document}
\maketitle

%----------------------------------------------------------------------
\begin{abstract}
Staff officers learning the Military Decision-Making Process (MDMP) need reliable
explanations of its steps, screening criteria, and war-gaming boundaries.
Many specialized assistants, however, are closed or depend on proprietary methods.
We ask whether public doctrine and fictional scenarios are sufficient to train an
open, publishable coaching model.
The training process excludes customer operational data and proprietary algorithms.
We then test whether the resulting model can beat a strong general-purpose baseline.
We construct a leak-reviewed instruction corpus, adapt
\texttt{Mistral-7B-Instruct-v0.3} with LoRA~\citep{hu2022lora}, and evaluate on
20 held-out golden questions.
The untuned base model passes $1/20$ ($5\%$).
Our best adapter passes $14/20$ ($70\%$).
The system is an unofficial educational prototype and is not a substitute for
authoritative Army doctrine~\citep{fm50,adp50}.
\end{abstract}

%----------------------------------------------------------------------
\section{Introduction}
\label{sec:intro}

\subsection{Motivation: MDMP coaching need vs closed/proprietary tooling}
\label{sec:intro-motivation}

MDMP structures how U.S.\ Army staffs turn missions into courses of action (COAs),
compare them, and produce orders~\citep{fm50,adp50}.
Learners and junior planners often need quick, consistent answers about which
step owns war gaming, when a decision matrix is used, and how FASDC screening
relates to later comparison.
Commercial or in-house assistants can help, but when they embed proprietary
decision algorithms or customer-specific planning data, they cannot be shared
as open research artifacts.
Other researchers also cannot reproduce or audit them.
An open coaching assistant must therefore be built from sources that are safe
to publish.
Its training text must not leak product names, customer identifiers, or closed
methods.

\subsection{Research question: Can open doctrine alone yield a competitive open assistant?}
\label{sec:intro-rq}

We study whether public doctrine summaries and a fictional training scenario
suffice to specialize a small open instruct model so that it outperforms the
same untuned model on held-out MDMP questions.
Open construction and leak controls need not preclude useful gains over a strong
general-purpose baseline.
We answer empirically with a LoRA adapter trained only on reviewed open pairs
and scored on a golden set that never appears verbatim in training.

\subsection{Contributions}
\label{sec:intro-contrib}
\begin{itemize}[leftmargin=*]
  \item An open MDMP instruction resource grounded in public doctrine markdown
        and a fictional scenario (OBJ FALCON), totaling 324 reviewed pairs after
        conflict repair and failure-targeted expansion.
  \item A leak-review gate that rejects proprietary product terms, customer-tied
        designations, and related closed\-method vocabulary before training.
  \item A reproducible LoRA fine-tuning recipe on
        \texttt{Mistral-7B-Instruct-v0.3}~\citep{mistral7b} (4-bit, Unsloth)
        suitable for a single modern GPU workstation.
  \item A 20-question golden evaluation showing the adapter improves from
        $1/20$ ($5\%$) on the base model to $14/20$ ($70\%$) on the best run.
\end{itemize}

\subsection{Paper roadmap}
\label{sec:intro-roadmap}

Section~\ref{sec:related} situates MDMP coaching, LoRA fine-tuning, and open-data
leakage risks.
Section~\ref{sec:method} describes publishability constraints, corpus and pair
construction, training, and golden scoring.
Section~\ref{sec:results} leads with base versus adapted performance, then
supporting iteration analyses.
Section~\ref{sec:discussion} discusses open release, limitations, and responsible use.

%----------------------------------------------------------------------
\section{Background and related work}
\label{sec:related}

\subsection{MDMP and staff-planning assistance}
\label{sec:related-mdmp}

MDMP comprises seven steps~\citep{fm50,adp50}:
\begin{enumerate}[leftmargin=*]
  \item Receipt of Mission
  \item Mission Analysis
  \item COA Development
  \item COA Analysis (war gaming)
  \item COA Comparison
  \item COA Approval
  \item Orders Production
\end{enumerate}
This paper uses doctrine as factual grounding for instruction tuning.

\subsection{Domain adaptation with LoRA}
\label{sec:related-lora}

Full fine-tuning of multi-billion-parameter models is costly.
Parameter-efficient fine-tuning (PEFT) adapts a pretrained model without
updating all of its weights.
Low-Rank Adaptation (LoRA) is one PEFT method: it freezes the base weights and
trains small adapter matrices~\citep{hu2022lora}, enabling domain specialization
with modest hardware.
We apply LoRA to an open 7B instruct model~\citep{mistral7b} using 4-bit
loading and Unsloth tooling as an engineering convenience, not as a claimed
algorithmic contribution.

\subsection{Open-weight models, public datasets, and leakage risks in specialized domains}
\label{sec:related-open}

Open weights and public datasets make specialized assistants easier to share.
However, domain corpora can accidentally include proprietary product names,
customer units, or closed algorithms.
Such leakage undermines both legal publishability and scientific
reproducibility.
Our pipeline therefore treats leak review as a first-class method step rather
than a post hoc checklist.

\subsection{Positioning: open doctrine coaching vs proprietary defense NLP}
\label{sec:related-position}

This work targets a publishable coaching assistant for MDMP vocabulary and
step boundaries, not a classified planning system or a proprietary decision
engine.
We restrict training data to open doctrine summaries and fiction.
We also filter vocabulary associated with closed methods.
These controls allow release without exposing customer or commercial IP.
The evaluation then measures whether this restricted approach beats the
untuned base model.

%----------------------------------------------------------------------
\section{Method: an open, publishable adaptation pipeline}
\label{sec:method}

\subsection{Publishability constraints}
\label{sec:method-constraints}

Allowed sources are (i)~summaries of publicly described MDMP concepts and
(ii)~fictional scenarios with no real unit or customer binding.
Forbidden content includes customer operational data, proprietary algorithm
or product terminology, and designations that would re-identify closed work.
A scripted leak review scans training JSONL for forbidden patterns
(Appendix~\ref{app:leak}) and must pass before each train split.
These constraints encode the requirement that the assistant remain publishable
and free of closed IP.

\subsection{Public corpus and fictional scenario design}
\label{sec:method-corpus}

The corpus includes five doctrine markdown files.
They cover MDMP steps, war gaming, COA screening, a glossary, and source
references.
The sources note points readers to \mbox{FM 5-0}, \mbox{ADP 5-0}, and other
public references listed in the repository~\citep{fm50,adp50}.
A single fictional scenario (OBJ FALCON) supplies scenario-coaching examples
without real-world capture data.
All instructional claims are treated as educational summaries, not
authoritative doctrine text.

\subsection{Instruction dataset construction and conflict harmonization}
\label{sec:method-data}

Pairs follow a JSONL schema with fields
\texttt{instruction}, \texttt{input}, \texttt{output}, \texttt{bucket},
\texttt{source}, and \texttt{reviewed}.
Buckets span MDMP steps, step boundaries, glossary, war gaming, COA screening,
and scenario coaching.
Reviewed pairs are split $85\%/15\%$ into train and held-out loss eval files.
The dataset first grew to roughly 300 pairs.
A subsequent conflict review removed self-contradictory Step~4/Step~5 guidance
and added harmonized replacements.
A later failure-targeted batch raised the corpus to 324 reviewed pairs.
After the final split, 275 pairs were used for training and 49 for loss
evaluation.
Golden questions are never copied verbatim into training files.
Conflict repair helps keep the open-only dataset internally consistent enough
for instruction tuning.

\subsection{Model and LoRA training protocol}
\label{sec:method-train}

We use \texttt{mistralai/Mistral-7B-Instruct-v0.3}~\citep{mistral7b} as the base
model.
The choice is deliberate: a 7B open instruct model supports fast LoRA iteration,
fits in roughly 12\,GB under 4-bit training, and keeps the debug surface small
while the data and evaluation pipeline are still being proven.
Larger open models remain available for later quality runs once that pipeline
is stable.
We do not claim that 7B is the strongest possible MDMP coach.

We fine-tune this base with LoRA rank $r{=}16$, $\alpha{=}32$, and dropout
$0.05$.
Training uses two epochs and a per-device batch size of 2.
Four gradient-accumulation steps produce an effective batch size of 8.
The learning rate is $2{\times}10^{-4}$, the sequence length is 2048, and the
model is loaded in 4-bit form.
Training ran on an NVIDIA GB10 workstation via Unsloth.
Prompts use the Mistral instruct format
\texttt{<s>[INST] \ldots{} [/INST]}.
Adapters are written to \texttt{outputs/mistral7b-mdmp-lora} for inference and
evaluation.
Inference needs less hardware than training: modest GPU memory for 4-bit
plus LoRA, or CPU after a merged quantized export
(Table~\ref{tab:hardware}).

\begin{table}[H]
  \centering
  \caption{Approximate hardware for training versus running the LoRA-adapted
  Mistral-7B assistant. Figures are order-of-magnitude guides for single-request
  use, not formal benchmarks.}
  \label{tab:hardware}
  \begin{tabular}{p{0.28\linewidth}p{0.22\linewidth}p{0.38\linewidth}}
    \toprule
    Setup & Approx.\ memory & Notes \\
    \midrule
    Training (4-bit LoRA / QLoRA) & ${\sim}$12\,GB+ GPU &
      Used on NVIDIA GB10; not the theoretical minimum \\
    GPU inference (4-bit + LoRA) & ${\sim}$5--8\,GB VRAM &
      Current Unsloth/transformers load path \\
    Comfortable consumer GPU & 8--12\,GB VRAM &
      Short coaching replies, batch size 1 \\
    Tight GPU & ${\sim}$6\,GB VRAM &
      Possible in 4-bit; watch context length \\
    CPU after merge + quantize & 8--16\,GB system RAM &
      e.g.\ GGUF / \texttt{llama.cpp}; slower than GPU \\
    \bottomrule
  \end{tabular}
\end{table}

\subsection{Held-out golden evaluation}
\label{sec:method-eval}

Product quality is measured with 20 golden questions maintained separately from
the train/eval JSONL split.
Each item specifies \texttt{must\_mention} and optional
\texttt{must\_not\_mention} terms.
A rule-based scorer marks a pass only when all required terms appear and no
forbidden terms appear.
Decode settings are held fixed across base and adapter runs
(temperature $0.1$, up to 256 new tokens).
This design tests whether the open-only adapter improves substantive MDMP
answers relative to the untuned open base model.

%----------------------------------------------------------------------
\section{Experiments and results}
\label{sec:results}

\subsection{Setup: base model as the open general-purpose baseline}
\label{sec:results-setup}

The primary comparison is the same open instruct model with and without the
LoRA adapter, using identical generation settings and the same 20 golden
questions.
The metric is golden pass rate (number of questions that satisfy the
must-mention rules).
Training-loss evaluation on the $15\%$ JSONL split is used only as a training
monitor, not as the product score.

\subsection{Main result: fine-tuned open assistant vs base}
\label{sec:results-main}

Table~\ref{tab:golden-main} reports the headline result.
The untuned base model passes $1/20$ ($5\%$).
Our best adapter (v7) passes $14/20$ ($70\%$).
That gap is the empirical evidence that an open-doctrine, leak-reviewed LoRA
adaptation can substantially outperform the general-purpose baseline on
held-out staff-planning questions.

\begin{table}[t]
  \centering
  \caption{Golden-question pass rate on 20 held-out MDMP questions.
  The base model has no adapter. The fine-tuned run is the best LoRA checkpoint (v7).}
  \label{tab:golden-main}
  \begin{tabular}{lcc}
    \toprule
    Model & Pass & Rate \\
    \midrule
    Base Mistral-7B-Instruct (no adapter) & $1/20$ & $5\%$ \\
    Fine-tuned LoRA (v7, best run) & $14/20$ & $70\%$ \\
    \bottomrule
  \end{tabular}
\end{table}

\subsection{Supporting analyses: data iteration, targeted repair, failure modes}
\label{sec:results-support}

Table~\ref{tab:golden-progress} summarizes intermediate adapters.
Early runs with few pairs underperform, while mid-scale runs improve.
A broad expansion to more pairs (v5) \emph{regressed} relative to v4
($40\%\rightarrow 35\%$).
Conflict harmonization (v6) and failure-targeted pairs (v7) recovered and
exceeded prior peaks ($55\%$ then $70\%$).
These supporting results do not replace the base-versus-adapted claim.
They show that open-only data must be curated.
Volume alone is not enough.

\begin{table}[t]
  \centering
  \caption{Golden pass rate across training iterations (same 20 questions).
  Train pairs are the JSONL split size used for that run.}
  \label{tab:golden-progress}
  \begin{tabular}{lccc}
    \toprule
    Run & Train pairs & Pass & Rate \\
    \midrule
    baseline & n/a & $1/20$ & $5\%$ \\
    v1 & 42 & $0/20$ & $0\%$ \\
    v2 & 153 & $5/20$ & $25\%$ \\
    v3 & 178 & $6/20$ & $30\%$ \\
    v4 & 203 & $8/20$ & $40\%$ \\
    v5 & 255 & $7/20$ & $35\%$ \\
    v6 & 258 & $11/20$ & $55\%$ \\
    v7 & 275 & $14/20$ & $70\%$ \\
    \bottomrule
  \end{tabular}
\end{table}

\subsection{What remains hard under an open-only regime}
\label{sec:results-hard}

Six golden items still fail under v7.
Some answers omit the literal phrase ``Step~$N$,'' such as ``5, COA
Comparison.''
Other answers omit FASDC purpose terms or assign criteria development and the
decision matrix to the wrong steps.
The scorer is also brittle when it requires ``synchronization'' but the answer
uses ``synchronize.''
Open doctrine alone therefore improves competence markedly but does not yet
guarantee fully robust step phrasing or terminology coverage on every probe.

%----------------------------------------------------------------------
\section{Discussion and conclusion}
\label{sec:discussion}

\subsection{Implications for open release}
\label{sec:discussion-release}

The natural release package has three parts: a LoRA adapter, the reviewed
instruction dataset, and a model card.
The model card should document license terms, data scope, leak-review policy,
and golden scores.
Training excludes customer data and proprietary algorithm vocabulary.
The artifacts are therefore suitable for public Hub distribution, subject to
the base model's license~\citep{mistral7b}.
They are also practical to try: Table~\ref{tab:hardware} summarizes modest GPU
requirements and CPU use after a merged, quantized export.
Official doctrine is cited by reference rather than reproduced wholesale.

\subsection{Limitations and responsible use}
\label{sec:discussion-limits}

The assistant is unofficial and educational.
It is not affiliated with the U.S.\ Army and must not replace
\mbox{FM 5-0} / \mbox{ADP 5-0}~\citep{fm50,adp50}.
The golden set is small ($n{=}20$), English-only, and U.S.\ MDMP-framed.
A high golden pass rate does not imply robustness to paraphrase:
informal probes with near-miss wording (for example asking for ``the steps of MDMP''
rather than the golden ``seven steps \ldots\ in order'', or ``What is step~5/7?'')
produced truncated lists and off-by-one step answers even when the corresponding
golden items passed.
A retrieval-augmented control that grounds generation in the same public
doctrine corpus was more stable on those factual lookups; that comparison was
not a full head-to-head benchmark and is not part of the published LoRA release.
We did not formally evaluate multi-turn staff coaching or classified scenarios.
Automatic scoring can be brittle to near-synonyms.
Users should verify any operationally consequential answer against authoritative
doctrine.

\subsection{Conclusion}
\label{sec:discussion-conclusion}

We adapted a 7B open instruct model with a leak-reviewed, open-doctrine corpus.
The LoRA adapter raised held-out MDMP golden accuracy from $5\%$ to $70\%$.
It did so without proprietary algorithms or customer data.
This result shows that a publishable coaching assistant is attainable under
those constraints.

%----------------------------------------------------------------------
\section*{Acknowledgments}
We thank colleagues who reviewed training pairs for doctrine consistency and
leak risk.
Any errors are the author's alone.
Opinions expressed are those of the author and do not represent official
positions of Decision Lens Inc.\ or the U.S.\ Army.

\bibliographystyle{plainnat}
\bibliography{refs}

%----------------------------------------------------------------------
\appendix

\section{Golden questions}
\label{app:golden}

The golden set probes step order, war gaming versus comparison, and FASDC
screening.
It also covers war-game methods, action-reaction-counteraction, evaluation
criteria timing, and the decision matrix.
The remaining questions address COA development, orders products, commander's
intent, COA approval, synchronization matrices, CCIR, mission analysis, and
COA definition.
Full must-mention lists appear in the repository file
\texttt{eval/golden\_questions.json}.

\section{Training hyperparameters}
\label{app:hparams}

\begin{table}[h]
  \centering
  \caption{LoRA fine-tuning configuration (from \texttt{train/config.yaml}).}
  \label{tab:hparams}
  \begin{tabular}{ll}
    \toprule
    Parameter & Value \\
    \midrule
    Base model & \texttt{mistralai/Mistral-7B-Instruct-v0.3} \\
    Max sequence length & 2048 \\
    Load in 4-bit & true \\
    LoRA $r$ / $\alpha$ / dropout & 16 / 32 / 0.05 \\
    Epochs & 2 \\
    Per-device batch size & 2 \\
    Gradient accumulation & 4 \\
    Learning rate & $2{\times}10^{-4}$ \\
    Warmup ratio & 0.03 \\
    Hardware & NVIDIA GB10 \\
    \bottomrule
  \end{tabular}
\end{table}

\section{Leak-review pattern families}
\label{app:leak}

The leak-review script rejects several high-level pattern families.
These include proprietary product or platform names tied to closed tooling and
customer or service abbreviations associated with capture work.
The script also rejects proprietary analytic vocabulary, such as closed scoring
or simulation jargon.
Finally, it rejects real or scenario designations reserved as non-publishable
relative to the open fictional set.
Exact regular expressions live in \texttt{scripts/leak\_review.py} and are
intentionally maintained outside the public narrative listing so that the
paper describes policy without serving as a bypass guide.

\end{document}
